
\section{Knorrer reduction}\label{sec: Knor}
\subsection{Construction of the Knorrer functor} In this section we introduce a slightly smaller space $\calXr_2$ which could be used to model our convolution algebra. The advantage of this perspective is that the generators of the
braid group have a simpler description.

Let $\calXr_2:=\frb\times G\times \frn$ and let us define the following potential on it
\begin{equation}\label{eq: red potential}
\SymbolPrint{\Wr}(X,g,Y)=\Tr(X\Ad_g(Y)).
\end{equation}

This potential is $B\times B$ equivariant with respect to the following action
$$(b_1,b_2)\cdot (X,g,Y)=(\Ad_{b_1}(X) , b_1 g b_2^{-1},\Ad_{b_2}(Y)).$$
The space $\calX_2$ projects  to the reduced space $\calXr_2$ by the projection $\bar{\pi}$:
$$\bar{\pi}(X,g_1,Y_1,g_2,Y_2)=(\Ad_{g_1}^{-1}(X)_+,g_1^{-1}g_2,Y_2).$$
The map is $B$-equivariant with respect to the second copy of $B$ in $B^2$.

The category of matrix factorizations $\MF_{B^2}(\calXr_2,\Wr)$ is equivalent to our main category
\begin{proposition}\label{prop: Knorrer} There is an embedding of categories
$$\SymbolPrint{\Phi}: \MF_{B^2}(\calXr_2,\Wr)\to \MF_{B^2}(\calX_2,W).$$
\end{proposition}
\begin{proof} Using  the elementary properties of the trace we obtain:
$$\Tr(\Ad^{-1}_{g_1}(X)Y_1-\Ad^{-1}_{g_{2}}(X)Y_{2})=\Tr(\tilde{X}_{1+}+\Ad_{g_{12}}(Y_2))-\Tr(\tilde{X}_{1--}(Y_{1}-\Ad_{g_{12}}(Y_2)_{++})),$$
where $\SymbolPrint{ g_{12}}=g_1^{-1}g_2,\quad \SymbolPrint{\tilde{X}_{i}}=\Ad^{-1}_{g_{i}}(X).$ Let us denote the first summand in the last formula as $\Wr$.

The last summand in the last formula is quadratic hence the non-equivariant version of the statement is exactly theorem 3.1 of \cite{Kn}.
The above mentioned theorem from \cite{Kn} provides the functor \(H: \textup{MF}(\CC^n,f)\to \textup{MF}(\CC^{n}\times\CC^2,f+vu)\). It is shown in \cite{Or} that
\(\CC^n\) in the theorem can be replaced by any affine variety and we use the result of \cite{Or} here since
\begin{equation}\label{eq:prod X}
  \calX_2=\calXr_2\times \CC^{2N}
  \end{equation}
with \(\tilde{X}_{1--}\) and \(Y_1-\Ad_{g_{12}}(Y_2)_{++}\) providing coordinates along \(\CC^{2N}\).
In particular, rewriting the functor
\(H\) in our notations  we obtain the functor $\Phi$:
\begin{equation}\label{eq: Knorrer functor}
 \Phi'(\mathcal{F})=\bar{\pi}^*(\mathcal{F})\otimes \begin{bmatrix} (\tilde{X}_{1--} )& (Y_{1}-\Ad_{g_{12}}(Y_2)_{++})\end{bmatrix},
 \end{equation}
where the last term is the Koszul matrix factorization and $(M)$ stands for the column composed of the non-trivial entries of $M$.

The equation (\ref{eq: Knorrer functor}) defines a functor $\Phi': \MF(\calXr_2,\Wr)\to \MF(\calX_2,W)$ which gives a non-equivariant version of our embedding.
Unfortunately, the map $\bar{\pi}$ is not $B^2$-equivariant hence we can not use the pull back along this map to define
an equivariant version of the functor $\Phi$. To get around this problem we give another interpretation of the functor $\Phi$ that extends to the equivariant setting.

In our construction we use the space $\SymbolPrint{\widetilde{\calX}_2}=\frb\times G\times \frn\times G\times \frn$ that projects on $\calXr_2$:
$$\SymbolPrint{\pi_y}: \widetilde{\calX}_2\to \calXr_2,\quad \pi_y(\Ad_{g_1}^{-1}(X)_+,g_1,Y_1,g_2,Y_2)=(X,g_1^{-1}g_2,Y_2),$$
and is isomorphic to the sub variety of  $\calX_2$  defined by the ideal $I_{--}$ generated by the entries of the matrix $\tilde{X}_{1--}$:
$$ \SymbolPrint{j^x}: \widetilde{\calX}_2\to \calX_2.$$

Since we have the product decomposition (\ref{eq:prod X}) and the entries of \(\tilde{X}_{1--}\)  are coordinates along \(\CC^{2N}\),
the generators of the ideal $I_{--}$ form a regular sequence and $\pi_y^*(\Wr)|_{\widetilde{\calX}_2}=W|_{\widetilde{\calX_2}}$.
Hence, according to the subsection~\ref{ssec: eq push f}, the functor \[j^x_*: \MF_{B^2}(\widetilde{\calX}_2,\pi_y^*(\Wr))\to \MF_{B^2}(\calX_2,W)\] is
well defined.
Now let us observe that $\Phi'(\calF)=j^x_*\circ \pi_y^*(\calF)^\sharp$. Indeed, $\bar{\pi}^*(\calF)\otimes [(\tilde{X}_{1,--}),0]$ is an element of 
$\textup{Com}_{per}^\ell(\calX_2)$ and the elements $\tilde{X}_{1,--}$ form a regular sequence, hence $\Phi'(\calF)$ is a unique extension of
$\bar{\pi}^*(\calF)\otimes [(\tilde{X}_{1,--}),0]$ to an element of $\MF(\calX_2,W)$ from the lemma~\ref{lem: ext MF+eq}.
On the other hand since $\bar{\pi}|_{\widetilde{\calX_2}}=\pi_y$, thus $\bar{\pi}^*(\calF)$ is an extension of $\pi_y^*(\calF)$ on $\calX_2$ (as required by 
the first step of our construction of $j_*^x$ from the  section~\ref{ssec: noneq push f}). Thus $j_*^x\circ\pi_y^*(\calF)$ is also an extension of
$\bar{\pi}^*(\calF)\otimes [\tilde{X}_{1,--},0]$ to an element of $\MF(\calX_2,W)$. Hence by the uniqueness result we have $\Phi'(\calF)=j_*^x\circ\pi_y^*(\calF)^\sharp$.
 


% $\Phi(\calF)|_{\widetilde{\calX}_2}=\pi_y^*(\calF)\otimes \mathrm{K}(I_{--})$ and
% $\Phi(\calF)$ provides and extension of the generators of the ideal to the element of $\MF(\calX_2,W)$. But so is $j^x_*\circ \pi_y^*(\calF)$ and the
% uniqueness of the extension from lemma~\ref{lem: uniq MF+eq} implies $\Phi(\calF)=j^x_*\circ\pi_y^*(\calF)$.


%Since the projection $\bar{\pi}$ is $B$ equivariant, the pull back $\bar{\pi}^*(\calF)$ acquires a natural $B$-equivariant structure with respect to the second copy of $B_{(2)}$ in
%$B^2$.
%The second factor in (\ref{eq: Knorrer functor}), which we call $CC$, has columns equivariant with respect to $B_{(2)}$ thus $\Phi(\calF)$ acquires a natural $B_{(2)}$-equivariant structure.

%On the other hand neither $\bar{\pi}^*(\calF)$ nor $CC$ have a natural $B_{(1)}$-equivariant structure because the map $\bar{\pi}$ is not $B_{(2)}$-equivariant and below we explain how one can introduce
%a $B_{(2)}$-equivariant structure on $\Phi(\calF)$. First let us us fix notation $\hat{Y}=Y_1-\Ad_{g_{12}}(Y_2)_{++}$ and $E_{ij}$, $1\le i<j\le n$ is the basis of $\frn$.
There is a unique $B^2$ equivariant structure on $\widetilde{\calX}_2$ that makes  the map $\pi_y$  $B^2$-equivariant. Let's fix this $B^2$-equivariant structure on $\widetilde{\calX}_2$.
Moreover, the ideal $I_{--}$ is also $B^2$-invariant. Let $\hat{Y}:=Y_{1}-\Ad_{g_{12}}(Y_2)$.
Then the variables $\hat{Y}_{++},  Y_2, \tilde{X}_1, g_{12}, g_1$ form a coordinate system on $\calX_2$ and are $B^{(2)}$-equivariant for trivial reasons.
On the other hand, the action of the Lie algebra of the unipotent part of $B^{(1)}$ is given
by
\[ E_{ij}\cdot \tilde{X}_{1,kl}=\delta_{jl} \tilde{X}_{1,ki}-\delta_{ki} \tilde{X}_{1,jl},\quad  E_{ij}\cdot \hat{Y}_{kl}=\delta_{jl} \Ad_{g_{12}}(Y_2)_{ki}-\delta_{ki} \Ad_{g_{12}}(Y_2)_{jl},\]
\[E_{ij}\cdot (g_{12})_{kl}=
\delta_{ki}(g_{12})_{jl},\quad E_{ij}\cdot (g_1)_{kl}=\delta_{ki}(g_1)_{kl}
\]
 and the action on the rest of coordinates is trivial. Thus $I_{--}$ is $\frn^2$ invariant.
  Moreover, the above mentioned $B^2$-action on $\widetilde{\calX}_2$ coincides with the standard $B^2$-action on $\calX_2$ restricted to
 on $B^2$-invariant subvariety $\widetilde{\calX}_2$.

 Thus we defined a functor:
 $$\SymbolPrint{\Phi}:\MF_{B^2}(\calXr_2,\Wr)\to \MF_{B^2}(\calX_2,W),\quad \Phi(\calF)=j^x_*\circ \pi_y^*(\calF).$$




%%%%%%%%%%%%%%%%%%%%%%%%%%%%%%%%%%%%%%%%%%%%%%%%%%%%%%%%%%%%%%%%%%%%%%%%%%%%%%%%%%%%%%%%%%%%%%%%%%%%%%%%%%%%%%%%%%


To show that the functor $\Phi$ is an embedding we define the functor $\Psi : \MF_{B^2}(\calX_2,W)\to \MF_{B^2}(\calXr_2,\Wr)$ such that $\Psi\circ\Phi= 1$.
Since the ideal $I_{--}$ generated by the matrix elements of $\tilde{X}_{1,--}$ is $B^2$ invariant we have well defined functor
$\SymbolPrint{Rest}: \MF_{B^2}(\calX_2,W)\to \MF_{B^2}(\widetilde{\calX_2}, \pi_y^*(\Wr))$ which is given by taking a quotient by the ideal $I_{--}$.
The push-forward along the projection $\pi_y: \calX_2\times \frn\to \calX_2$ provides the last step of the construction for
$\Psi:=\pi_{y*}\circ Rest$.
\end{proof}

\subsection{Properties of the Knorrer functor} As in the previous proof, $CC$ denotes the correction complex $$CC:= \begin{bmatrix} (\tilde{X}_{1--} )& (Y_{1}-\Ad_{g_{12}}(Y_2)_{++})\end{bmatrix}$$
Let us explain the equivariant structure of this matrix factorization. Since $B_{(2)}$-action does not change entries of the complex $\SymbolPrint{CC}$, the $B_{(2)}$ equivariant structure of $B_{(2)}$ is trivial.

To describe $B_{(1)}$ equivariant structure, let us write the differential of $CC$ as
\[ \SymbolPrint{D^{kn}}=D^{kn}_++D^{kn}_-, \quad D^{kn}_+=\sum_{k>l}(\tilde{X}_{1,--})_{kl}\theta_{kl},\]
\[D^{kn}_-=\sum_{k>l}
%\hat{Y}_{lk}
 (Y_{1}-\Ad_{g_{12}}(Y_2)_{++})_{lk}
\frac{\partial}{\partial \theta_{kl}},\]
where  $\theta_{ij}$ are odd variables. The equivariant structure is defined by the action on $\theta_{kl}$:
$$ E_{ij}\cdot \theta_{kl}=\sum_{k>i}\theta_{ki}-\sum_{j>l}\theta_{jl}$$

The differential $D^{kn}$ is not $\frn_{(2)}$ invariant:
$$ E_{ij}\cdot D^{kn}_+=0,\quad E_{ij}\cdot D^{kn}_-=\sum_{j>l} [\Ad_{g_{12}}(Y_2)_-]_{il}\frac{\partial}{\partial\theta_{jl}}-\sum_{k>i} [\Ad_{g_{12}}(Y_2)_-]_{ik}\frac{\partial}{\partial\theta_{ki}}.$$
Thus the correction differentials are non-trivial in $\Phi(\calF)$ but the correction differentials are not arbitrary and below we emphasize properties of the correction
differentials that are needed for the proofs that follow.

\def\mltV{V_{\mathrm{m}}}
\begin{proposition} Let $\calF=(M,D,\partial_l,\partial_r)\in \MF_{B^2}(\calXr_2,\Wr)$, $M=\mltV\otimes \CC[\calXr_2]$, $\mltV\in \Mod^{filt}_{\frn^2}$
The matrix factorization $\Phi(\calF)\in \MF_{B^2}(\calX_2,W)$ is the equivariant matrix factorization
$$ (\mltV\otimes \CC[\calX_2], D+D^{kn}, \tilde{\partial}_l,\tilde{\partial}_r)$$
where $\tilde{\partial}_r=\partial_r$.
\end{proposition}
\begin{proof}
The only part that requires a discussion is the fact that the differential $\Phi(\calF)^{\sharp}$ has differential $D+D^{kn}$: a priori our construction of the push forward
$j^x_*(\pi_y^*(\calF))$ relies on the extension lemma~\ref{lem: ext MF+} and involves choice of an extension $\tilde{D}$ of $D$ from variety $\widetilde{\calX}_2$ to $\calX_2$
and completion of $\tilde{D}+D^{kn}_+$  to the differential $\tilde{D}+D^{kn}_++d^-$ such that $(\tilde{D}+D^{kn}_++d^-)^2=W$.

However, in our case map $j^x$ is actually a section
of the (non-equivariant) map: $$\pi_x:\calX_2\to \widetilde{\calX}_2,\quad \pi_x(X,g_1,Y_1,g_2,Y_2)=(\Ad_{g_1}^{-1}(X)_+,g_1^{-1}g_2,Y_2).$$
 Thus according to the remark~\ref{rem: section}
we can choose $\tilde{D}$ to be $D=\bar{\pi}^*(D)=\pi_x^*(\pi_y^*(D))$. Moreover, $d^-=D_-^{kn}$ because $(D+D^{kn})^2=W$ and the completion of  the differential is unique up to an isomorphism.
\end{proof}

% and matrix coefficients of $\tilde{\partial}_l$ are polynomials of variables $Y_2,X,g_1,g_2$ but not of
%$Y_1$.
%\en{prop}
%\begin{proof}
%The only statement that needs a discuss is the statement of about $\tilde{\partial}_l$. First let us notice that the property holds for $\partial_l$.
%The correction differential $\tilde{\partial}_l$ is constructed by the iterative procedure explain in the proof of lemma~\ref{lem: ext MF+eq}. On $m$-th step of the
%procedure we obtain a correction $\tilde{\partial}^{{m)}$. The Knorrer complex is the two periodic folding of the Koszul matrix factorization
%and iterative procedure uses degree cohomological degree of this complex (see proof of lemma~\ref{lem: ext MF+eq} for details). In particular $\tilde{\partial}\in Hom
 %The next  step of the procedure is the extraction of the leading term with respect to
%Koszul homological grading of the expression
%$$ (D+D^{kn}+d_{ce}+\partial_l+\partial_l+\tilde{\partial}_l^{(m)})^2-W.$$
%However

%\end{proof}

\subsection{Convolution on the reduced space}
In order to define the convolution on the smaller space $\calXr_2$, consider
%It is convenient to work with the smaller space $\calXr_2$ and in particular we can introduce the following version of convolution.
%Define
the space $\SymbolPrint{\calXr_3}:=\frb\times G^2\times\frn$ with the following $B^3$-action
$$(b_1,b_2,b_3)\cdot (X,g_{12},g_{23},Y)=(\Ad_{b_1}(X),b_1 g_{12} b_2^{-1}, b_2 g_{23} b_3^{-1}, \Ad_{b_3}(Y)).$$
 There are the following
%$B^3$-equivariant
maps $\SymbolPrint{\bar{\pi}_{ij}}:\calXr_3\to\calXr_2$:
\[ \bar{\pi}_{12}(X,g_{12},g_{13},Y)=(X,g_{12},\Ad_{g_{23}}(Y)_{++}),\] \[\bar{\pi}_{23}(X,g_{12},g_{13},Y)=(\Ad_{g_{12}}^{-1}(X)_+,g_{23},Y),\]
\[\bar{\pi}_{13}(X,g_{12},g_{13},Y)=(X,g_{12}g_{23},Y).\]
%To ensure that all the maps are $B^3$ equivariant we trivially extend the $B^2$-action on $\calXr_2$ to $B^3$-action.


\begin{proposition}
$\bar{\pi}_{12}^*(\Wr)+\bar{\pi}_{23}^*(\Wr)=\bar{\pi}_{13}^*(\Wr)$.
\end{proposition}
\begin{proof}
Indeed we have
\begin{multline*}
\bar{\pi}_{12}^*(\Wr)+\bar{\pi}_{23}^*(\Wr)=\Tr(X\Ad_{g_{12}}(\Ad_{g_{23}}(Y)_{++}))+\Tr(\Ad_{g_{12}}^{-1}(X)_+\Ad_{g_{23}}(Y))=\\
\Tr(\Ad^{-1}_{g_{12}}(X)\Ad_{g_{23}}(Y)_{++})+\Tr(\Ad^{-1}_{g_{12}}(X)_+\Ad_{g_{23}}(Y)).
\end{multline*}
Thus the statement follows is implied by
\begin{multline*} \Tr(A B_{++})+\Tr(A_+B)=\Tr(AB)-\Tr(AB_-)+\Tr(A_+B_-)\\
  =\Tr(AB)+\Tr(A_-B_-)=\Tr(AB).
\end{multline*}

\end{proof}

For given $\calF,\calG\in \MF_{B\times B}(\calXr_2,W)$ it is tempting  to define
\begin{equation}\label{eq: red convolution prelim}
\calF\SymbolPrint{\bar{\star}}\calG:=\bar{\pi}_{13*}(\CE_{\frn^{(2)}}( \bar{\pi}_{12}^*(\calF)\otimes\bar{\pi}_{23}^*(\calG))))^{T^{(2)}},
\end{equation}
where $\frn^{(2)}$, $T^{(2)}$ are the subalgebras of the middle factor in $\frb^3$ acting on $\calXr_3$. The only issue with this definition is that
the maps $\bar{\pi}_{12}$ and $\bar{\pi}_{23}$ are not $B^3$ equivariant. Thus we have to define the $B^3$-equivariant structure on
the pull back $\bar{\pi}_{12}^*(\calF)\otimes\bar{\pi}_{23}^*(\calG)$ and we do it below.



\subsection{Reduced vs. non-reduced convolution}
%Let us introduce the following correcting $B^3$-equivariant complex:
%$$ CC_{123}=[\tilde{X}_{1,--}, Y_1-\Ad_{g_{12}}(\Ad_{g_{23}}(Y_3)_{++})_{++}].$$

%\begin{proposition}
%$$CC_{123}\in \MF_{B^3}(\calX_3,\pi_{13}(W-\bar{\pi}(W^{red})).$$
%\end{proposition}
%\begin{proof}
%It is to show the statement on an open subset of $\calX_3$. In particular we can choose a subset $\calX_3^{\circ}$ defined by conditions
%$$g_{12}=g_1^{-1}g_2\in B_-\cdot B_+,\quad g_{23}=g_2^{-1}g_3\in B_-\cdot B_+.$$
%On this set we can use $B^3$-equvariance of the complex and reduce our statement to the case
%$$g_{12}\in B_-,\quad g_{23}\in B_-.$$
%Finally under this assumptions we have
%$$\Ad_{g_{12}}(\Ad_{g_{23}}(Y_1)_{++})_{++}=\Ad_{g_{13}}(Y_1)_{++}$$
%and hence the complex $CC_{123}$ coincides with the correction complex $\pi_{13}^*(CC)$ from the functor
%$\Phi$ and the statement follows.
%\end{proof}

%\begin{proposition} For any $\calH,\calG\in \MF_{B\times B}(\calXr_2,W)$ we  have
%$$\bar{\pi}^*(\calF\bar{\star}\calG)\otimes \mathrm{H}^*(\frn^{(2)},CC_{123})=\Phi(\calF)\star\Phi(\calG).$$
%\end{proposition}

Let $\bar{\pi}_{con}:\calX_3\to \calXr_3$ be the map defined by
$$\SymbolPrint{\bar{\pi}_{con}}(X,g_1,Y_1,g_2,Y_2,g_3,Y_3)=(\Ad_{g_1}^{-1}(X)_+,g_1^{-1}g_2,g_2^{-1}g_3,Y_3).$$

\begin{proposition} \label{prop: Knorrer factor} Suppose $\calF,\calG\in \MF_{B^2}(\calXr_2,\Wr)$ then there is a sequence of  row transformations that identifies
$\pi_{12}^*(\Phi(\calF))\otimes_B\pi_{23}^*(\Phi(\calG))$ with the $B^3$-equivariant matrix factorization $\calH=(M,D,\partial_l,\partial_m,\partial_r)$,
\begin{equation}\label{eq: cor diffs}
\partial_l,\partial_m,\partial_r:
\Lambda^\bullet \frn \otimes M_{\star}\to \Lambda^{<\bullet}\otimes M_{\ge \star},
\end{equation}
$$\calH^\sharp=(M,D)=\bar{\pi}^*_{con}( \bar{\pi}_{12}^*(\calF))\otimes \bar{\pi}^*_{con}(\bar{\pi}_{23}^*(\calG))\otimes  [Y_2 -\Ad_{g_{23}}(Y_3)_{++},0]\otimes\pi_{13}^*(CC)$$
\end{proposition}


This proposition allows us to construct $B^3$-equivariant structure on $\bar{\pi}_{12}^*(\calF)\otimes\bar{\pi}_{23}^*(\calG)$.
First, we introduce few auxiliary spaces and map: $\SymbolPrint{\widehat{\calX}_3}:=\frg\times \frn\times G^3\times \frn$ and $\hat{\pi}:\calX_3\to\widehat{\calX}_3$ is given by
$$\SymbolPrint{\hat{\pi}}(X,g_1,Y_2,g_2,Y_2,g_3,Y_3)=(X,Y_1,g_1,g_2,g_3,Y_3).$$
We can embed $\widehat{\calX}_3$ into $\calX_3$ by
$$ \SymbolPrint{\hat{j}}:\widehat{\calX}_3\to \calX_3,\quad\hat{ j}(X,Y_1,g_1,g_2,g_{3},Y_3)\mapsto (X,g_{1},Y_1,g_2,\Ad_{g_{23}}(Y_3)_{++},g_{3},Y_3).$$
The map $\hat{j}$ the section of $\hat{\pi}$.

On the other hand, the matrix factorization $\calH^\sharp$ has Koszul complex $ [Y_2 -\Ad_{g_{23}}(Y_3)_{++},0]$ as one of the factors.
In other words, $\calH^\sharp=(M,D''+D')$, where $D'=\sum_{ij} (Y_2 -\Ad_{g_{23}}(Y_3)_{++})_{ij}\frac{\partial}{\partial\theta_{ij}}$.
Since the differential has the property (\ref{eq: cor diffs}) we contract $\calH$ along the differential $D'$ to obtain a homotopy  equivalence of
the matrix factorizations over $\CC[\widehat{\calX}_3]$
$$\calH\sim \widehat{\calH}=(\tilde{M},\tilde{D},\tilde{\partial}_l,\tilde{\partial}_m,\tilde{\partial}_r)\in \MF_{B^2}(\widehat{\calX}_3,\hat{j}^*\circ\pi_{13}^*(W))$$
where $\tilde{D}=\hat{j}^*(D'')$, $\hat{\partial}_l=\hat{j}(\partial_l)$, $\hat{\partial}_m=\hat{j}^*(\partial_m)$, $\hat{\partial}_r=\hat{j}^*(\partial_r)$.

Next, we introduce an analog of the Knorrer functor:
\[
  \SymbolPrint{\Phi^{(13)}}: \MF_{B^2}(\calXr_3,\bar{\pi}_{13}^*(\Wr))\to \MF_{B^2}(\widehat{\calX}_3,\hat{j}^*\circ\pi_{13}^*(W)).
\]
For that we introduce the space $$\SymbolPrint{\widetilde{\calX}_3}=\frb\times \frn \times G^3\times \frn$$ and fix a $B^3$-equivariant embedding
$\SymbolPrint{\tilde{j^x}}:\widetilde{\calX}_3\to \widehat{\calX}_3$, by identifying $\widetilde{\calX}_3$ with the subvariety of $\widehat{\calX}_3$ defined by the ideal
$\tilde{I}_{--}$ generated by the entries of $\Ad_{g_1}^{-1}(X)_{--}$.

The projection $\SymbolPrint{\tilde{\pi}_y}: \widetilde{\calX}_3\to \calXr_3$ defined by $$\tilde{\pi}_y(X,Y_1,g_1,g_2,g_3,Y_3)=(X,g_1^{-1}g_2,g_2^{-1}g_3,Y_3).$$
is $B^2$-equivariant. Thus since $\tilde{I}_{--}$ is $B^3$-equivariant we can define the functor:
$$ \Phi^{(13)}: \MF_{B^3}(\calXr_3,\Wr)\to  \MF_{B^3}(\widehat{\calX}_3,\hat{j}^*\circ\pi^*_{13}(W)),\quad \Phi(\calF)=\tilde{j}^x_*\circ\tilde{\pi}_y^*(\calF).$$

The following statement is analogous to the proposition~\ref{prop: Knorrer} and omit its proof:

\begin{proposition} The functor $\Phi^{(13)}$ provides an embedding of categories and
$$\Phi^{(13)}(\calF)^\sharp=\tilde{\pi}_{con}^*(\calF)\otimes [(\tilde{X}_{1,--}) (Y_1-\Ad_{g_{13}}(Y_3)_{++})],$$
where $\tilde{\pi}_{con}(X,Y_1,g_1,g_2,g_3,Y_3)=(\Ad_{g_1}^{-1}(X)_+,g_1^{-1}g_2,g_2^{-1}g_3,Y_3).$
\end{proposition}

Let us define a matrix factorization $\overline{\calH}\in \MF_{B^2}(\calXr_3,\Wr)$ in two steps. First, let us observe that there is
a well-defined functor $Rest: \MF_{B^2}(\widehat{\calX}_3,\hat{j}^*\circ\pi_{13}^*(W))\to \MF(\widetilde{\calX}_3,\tilde{\pi}_y^*(\Wr))$
which is the functor of taking quotient by the $B^3$-equivariant  ideal $\tilde{I}_{--}$. Second as a second step we apply
push forward along the projection map:
$$ \overline{\calH}:=\tilde{\pi}_{y*}\circ Rest(\widehat{\calH}).$$

The main result of this subsection is the following
\begin{proposition}
$$\widehat{\calH}=\Phi^{(13)}(\overline{\calH}).$$
\end{proposition}
\begin{proof}
Indeed, proposition~\ref{prop: Knorrer factor} implies that
$$\widehat{\calH}^\sharp=\tilde{\pi}^*_{con}( \bar{\pi}_{12}^*(\calF)\otimes \bar{\pi}_{23}^*(\calG))\otimes [(\tilde{X}_{1,--}) (Y_1-\Ad_{g_{13}}(Y_3)_{++})].$$
Hence $\overline{\calH}=\bar{\pi}_{12}^*(\calF)\otimes \bar{\pi}_{23}^*(\calG)$ and both
$\widehat{\calH}$ and $\Phi^{(13)}(\overline{\calH})$ provide an $B^2$-equivariant extension of $\tilde{\pi}_y^*(\overline{\calH})$ from $\widetilde{\calX}_3$ to
$\widehat{\calX}_3$. Hence the result follows from the uniqueness of the extension.
\end{proof}

The proof of the previous proposition explains that $\overline{\calH}$ is $B^3$-equivariant matrix factorization such that
$\overline{\calH}^\sharp=\bar{\pi}_{12}^*(\calF)\otimes\bar{\pi}_{23}^*(\calG).$ Thus we equipped the pull back
$\bar{\pi}_{12}^*(\calF)\otimes\bar{\pi}_{23}^*(\calG)$ with $B^3$-equivariant structure. We use notation
$$\SymbolPrint{(\bar{\pi}_{12}\otimes_B\bar{\pi}_{23})^*}(\calF\boxtimes\calG)\in \MF_{B^3}(\calX_3,\bar{\pi}_{13}^*(\Wr)).$$

We define binary operation on $\MF_{B^2}(\calXr_2,\Wr)$ by
$$\calF\bar{\star}\calG:=\bar{\pi}_{13 *}\CE_{\frn^{(2)}}\left((\bar{\pi}_{12}\otimes_B\bar{\pi}_{23})^*(\calF\boxtimes\calG)\right)^{T^{(2)}}.$$

The associativity of this binary operation follows from
\begin{corollary} For any $\calF,\calG\in\MF_{B^2}(\calXr_2,\Wr)$ we have
$$\Phi(\calF)\star\Phi(\calG)=\Phi(\calF\bar{\star}\calG).$$
\end{corollary}
\begin{proof}
There is a natural projection map $\hat{\pi}_{13}: \widehat{\calX}_3\to \calX_2$, $$\hat{\pi}_{13}(X,Y_1,g_1,g_2,g_3,Y_3)=(X,Y_1,g_1,g_3,Y_3).$$
Since we have homotopy $\hat{\calH}\sim \calH$, we conclude that
$$\Phi(\calF)\star\Phi(\calG)=\pi_{13*}(\calH)=\hat{\pi}_{13*}(\widehat{\calH})=\hat{\pi}_{13*}(\Phi^{(13)}(\overline{\calH})).$$
Thus to prove the statement we need to show that $\hat{\pi}_{13*}\circ \Phi^{(13)}=\Phi\circ\bar{\pi}_{13*}$ or equivalently
$\hat{\pi}_{13*}\circ \tilde{j}^x_*\circ \tilde{\pi}^*_y=j_*^x\circ\tilde{\pi}^*_y\circ\bar{\pi}_{13*}.$
The later identity holds because the map $\hat{\pi}_{13}\times\bar{\pi}_{13}: \widehat{\calX}_3\times\calXr_3\to \calX_2\times \widetilde{\calX}_2$ sends
the kernel of $\Phi^{(13)}$, which is $\tilde{j}^x\times \tilde{\pi}_y(\widetilde{\calX}_3)\subset  \widehat{\calX}_3\times\calXr_3$ to
the kernel of the map $\Phi$, which is $j^x\times \tilde{\pi}_y(\widetilde{\calX}_2)\subset  \calX_2\times\widetilde{\calX}_2$. More formally, we can proceed with the sequence of base isomorphisms as follows. We have the following diagram of maps:
\[
\begin{tikzcd}
\widehat{\calX_3}\arrow[d,"\hat{\pi}_{13}"]&\widetilde{\calX}_3\arrow[r,"\tilde{\pi}_y"]\arrow[d,dashed,"\tilde{\pi}_{13}"]\arrow[l,"\tilde{j}^x"]
&\overline{\calX}_3\arrow[d,"\bar{\pi}_{13}"]\\
\calX_2&\widetilde{\calX}_2\arrow[r,"\tilde{\pi}_y"]\arrow[l,"j^x"]&\overline{\calX}_2
\end{tikzcd}
\]
where the dashed arrow map $\tilde{\pi}_{13}$ is a unique smooth map such that the diagram commutes.
Then we can use change of base isomorphism to show:
\[
j_*^x\circ\tilde{\pi}_y^*\circ \bar{\pi}_{13*}=j_*^x\circ \tilde{\pi}_{13*}\circ\tilde{\pi}^*_y=\hat{\pi}_{13*}\circ\tilde{j}^x_*\circ\tilde{\pi}_y^*
\]


\end{proof}

Before explaining the proof of the proposition~\ref{prop: Knorrer factor} let us discuss of the elementary transformations in the context of our extension lemma~\ref{lem: ext MF+eq}.


\subsection{Row operations for generalized equivariant Koszul matrix factorizations}\label{ssec: row ops}
We use notations of the lemma~\ref{lem: ext MF+eq}
here and we also assume that the complex $(\calF_\bullet,d^+)$ is the Koszul complex tensored with matrix factorization $(M,D,\partial)$:
$$\calF_\bullet=\tilde{M} \otimes \Lambda^\bullet V, \quad d^+=\sum_{j=1}^n f_i\theta_i,$$
where $\theta_i$ form a basis of $V$ and $f_1,\dots,f_n$ form a regular sequence in $\CC[\calZ]$.
We call such matrix factorizations generalized equivariant Koszul matrix factorizations.

In these notations the corrective differentials that are computed by the iterative procedure from lemma~\ref{lem: ext MF+eq} are of the form
$$ d^-=\sum_{\vec{i}} d^-_{\vec{i}}\frac{\partial}{\partial \theta_{\vec{i}}}, \quad \partial^-=\sum_{\vec{i}}\partial^-_{\vec{i}} \frac{\partial}{\partial \theta_{\vec{i}}},$$
where $\theta_{\vec{i}}=\theta_{i_1}\dots \theta_{i_k}$ and $d^-_{\vec{i}}\in Hom_{\CC[\calZ]}(\tilde{M},\tilde{M})$,
$\partial_{\vec{i}}^-\in Hom(\Lambda^\bullet \frn,\Lambda^{<\bullet}\frn)\otimes Hom_{\CC[\calZ]}(\tilde{M},\tilde{M})$.

Let us use notation $D_{tot}:=d^++D+\partial+d^-+\partial^-+d_{ce}$.
There are three types of (generalized) row transformations for the complex $(\calF_\bullet,d^++D+d^-,\partial+\partial^-)$. The first type is given by automorphism
 $\exp(\phi\frac{\partial}{\partial\theta_k})$, $\phi\in \Hom_{\CC[\calZ]}(\calF,\calF)$ is an odd automorphism and the affect of this automorphism on the differential is:
 \begin{multline}
  \exp(-\phi\frac{\partial}{\partial \theta_k})\circ D_{tot}\circ \exp(\phi\frac{\partial}{\partial \theta_k})= \left(D+d_-+d_++[\phi,D]_+\frac{\partial}{\partial \theta_k}+f_k\phi\right)\\
 +d_{ce}+\partial+\partial^-+[\phi,d_{ce}+\partial+\partial_-]_+\frac{\partial}{\partial\theta_k}.
\end{multline}
That is the result of the row transformation of the first type is the equivariant matrix factorization
$$ (\tilde{M},D+d_-+d_++[\phi,D]_+\frac{\partial}{\partial \theta_k}+f_k\phi,\partial+\partial^-+[\phi,d_{ce}+\partial+\partial_-]_+\frac{\partial}{\partial\theta_k}).$$

The second type elementary row transformation is given by conjugation by the automorphism $\exp(\phi\theta)$. The result of conjugation of $D_{tot}$ by this automorphims is the following expression
$$ \exp(-\phi\theta_k)\circ D_{tot}\circ\exp(\phi\theta_k)=D_{tot}+D'+\partial'.$$
\begin{equation}\label{eq: sec row 1}
D'=[\phi,D]_+\theta_k+\sum_{k\in \vec{i} }d^-_{\vec{i}}\phi\frac{\partial}{\partial\widehat{\theta}_{\vec{i}}}+\theta_k\sum_{\vec{i}}[\phi,d^-_{\vec{i}}]\frac{\partial}{\partial\theta_{\vec{i}}},
\end{equation}
\begin{equation}\label{eq: sec row 2}
\partial'=\sum_{k\in \vec{i} }\partial^-_{\vec{i}}\phi\frac{\partial}{\partial\widehat{\theta}_{\vec{i}}}+\phi d_{ce}^k+
\theta_k(\sum_i[\phi,d_{ce}^i]\frac{\partial}{\partial\theta_i}+[\phi,\partial]_+ +\sum_{\vec{i}}[\phi,\partial^-_{\vec{i}}]\frac{\partial}{\partial\theta_{\vec{i}}}),
\end{equation}
where $\partial\widehat{\theta}_{\vec{i}}=\partial \theta_{i_1}\dots \widehat{\partial\theta_k}\dots\partial_{i_m}$, $\vec{i}=(i_1,\dots,i_m)$ and
$d_{ce}=\sum_i d_{ce}^i\frac{\partial}{\partial \theta_i}.$

The computation above indicates that if last summands in (\ref{eq: sec row 1}) and (\ref{eq: sec row 2}) vanish
\begin{equation}\label{eq: adm row tr}
\sum_{\vec{i}}[\phi,d^-_{\vec{i}}]\frac{\partial}{\partial\theta_{\vec{i}}}=0,\quad \sum_i[\phi,d_{ce}^i]\frac{\partial}{\partial\theta_i}+[\phi,\partial]_+ +\sum_{\vec{i}}[\phi,\partial^-_{\vec{i}}]\frac{\partial}{\partial\theta_{\vec{i}}}=0
\end{equation}
 then result of the row transformation is the equivariant matrix factorization:
 $$(\tilde{M},d^++D+d^-+[\phi,D]_+\theta_k+\sum_{k\in \vec{i} }d^-_{\vec{i}}\phi\frac{\partial}{\partial\widehat{\theta}_{\vec{i}}},\partial+\partial^-+\sum_{k\in \vec{i} }\partial^-_{\vec{i}}\phi\frac{\partial}{\partial\widehat{\theta}_{\vec{i}}}+\phi d_{ce}^k)$$

Finally the third type of row operation is given by automorphism $\exp(c\theta_i\frac{\partial}{\partial\theta_j})$ where $c\in \Hom_{\CC[\calZ]}(\calF,\calF)$ is an even morphism. This type of row transformation
was discussed in the subsection~\ref{ssec: str eq Kosz} for strongly equivariant matrix factorizations.
 We will use these type row operations only for strongly equivariant matrix factorization and omit a discussion
about complications that arise for weakly  equivariant generalized Koszul complexes.


The key feature of the row transformations discussed here is that differentials that encode the equivariant structure of the matrix factorization do not
contribute to the change of the non-equivariant shadow of the matrix factorization. Other words, if $T$ is one the transformations discussed above then
$$ T(\calF)^\sharp=T(\calF^\sharp).$$


\subsection{Proof of proposition~\ref{prop: Knorrer factor}}
%\begin{proposition} \label{prop: Knorrer homo}For any $\calH,\calG\in \MF_{B\times B}(\calXr_2,\Wr)$ we  have
%$$\Phi(\calF\bar{\star}\calG)=\Phi(\calF)\star\Phi(\calG),$$
%other words $\Phi$ is the convolution algebra homomorphism.
%\end{proposition}
%\begin{proof}

 %By the last remark from the previous subsection to prove statement it is enough to prove it for the non-equivariant

Let us expand the definition of convolution to obtain the following expression for $\pi_{12}^*(\Phi(\calF)^\sharp)\otimes\pi_{23}^*(\Phi(\calG)^\sharp)$:
\begin{multline}\label{eq: long MF}
\pi_{12}^*(\bar{\pi}^*(\calF^\sharp)\otimes CC)\otimes\pi_{23}^*(\bar{\pi}^*(\calG^\sharp)\otimes CC)\\
=(\pi_{12}\circ\bar{\pi})^*(\calF^\sharp)\otimes (\pi_{23}\circ\bar{\pi})^*(\calG^\sharp)\otimes \pi_{12}^*(CC)\otimes\pi^*_{23}(CC).
\end{multline}
The last factor is the Koszul complex $\begin{bmatrix} (\tilde{X}^2_{--} )& (Y_{2}-\Ad_{g_{23}}(Y_3)_{++})\end{bmatrix}$ hence we can use
row transformations of second type to eliminate variables $Y_2$  from the matrix factorization (\ref{eq: long MF}).
More precisely, use the row transformations of the second type to eliminate $Y_2$ from
the first factor $(\pi_{12}\circ \bar{\pi})^*(\calF)$  and from the second to last factors $\pi_{12}^*(CC)$.

 We can apply the second type of row transformation because
 $(\pi_{12}\circ \bar{\pi})^*(\calF^\sharp)$ and $\pi_{12}^*(CC)$ are the factors of $\pi^*_{12}(\Phi(\calF)^\sharp)$, thus differentials in these matrix factorizations
 commute with the differentials from the factor $\pi^*_{23}(\Phi(\calG)^\sharp)$ and hence the admissibility condition (\ref{eq: adm row tr}) is satisfied.
 On the other hand,  the second factor $(\pi_{23}\circ\bar{\pi})^*(\calG^\sharp)$ has no dependence on $Y_2$.

The result of elimination of $Y_2$ from the factor before the last  is $$CC_{123}=[\tilde{X}_{1,--}, Y_1-\Ad_{g_{12}}(\Ad_{g_{23}}(Y_3)_{++})_{++}]$$ and the first two factors are
\begin{gather*}
\pi^*_{12}(\bar{\pi}^*(\calF))|_{Y_2=\Ad_{g_{23}}(Y_3)_{++}}=\calF|_{Y=\Ad_{g_{23}}(Y_3)_{++},X_+=\tilde{X}_{1,+}}=\bar{\pi}^*_{con}(\bar{\pi}_{12}^*(\calF))\\
\pi^*_{23}(\bar{\pi}^*(\calG))|_{Y_2=\Ad_{g_{23}}(Y_3)_{++}}=\calG|_{Y=Y_3,X_+=\Ad_{g_{12}}^{-1}(\tilde{X}_{1})_+},
\end{gather*}
where $(X_+,g_{12},g_{23},Y)$ are coordinates on $\calXr_3$.
After the elimination the only term that has dependence on $Y_2$ is the Koszul complex  $\begin{bmatrix} G& (Y_{2}-\Ad_{g_{23}}(Y_3)_{++})\end{bmatrix}$ and since $\pi_{13}^*(W)$ does not depend
on $Y_2$ and the correction correction differentials $\partial$, which could depend on $Y_2$ does not contribute the square of the total differential, the element $G$ must vanish.


As the next step we use the complex $CC_{123}$ to eliminate dependence on $\tilde{X}_{1,--}$ from the other terms in the tensor product (\ref{eq: long MF}). For that
we use the first type of row transformation and we do not need to worry about the admissibility of our eliminations.
The term $\bar{\pi}^*_{con}(\bar{\pi}_{12}^*(\calF))$ has no dependence on $\tilde{X}_{1,--}$ so it is not affected by this elimination. On the other hand the second
term becomes:
$$ \calG|_{Y=Y_3,X_+=\Ad_{g_{12}}^{-1}(\tilde{X}_{1})_+, \tilde{X}_{1,--}=0}= \calG|_{Y=Y_3,X_+=\Ad_{g_{12}}^{-1}(\tilde{X}_{1,+})_+}=\bar{\pi}^*(\bar{\pi}_{23}^*(\calG)).$$

After this round of elimination the only dependence on $\tilde{X}_{1,--}$ in (\ref{eq: long MF}) comes from the Koszul complex $[\tilde{X}_{1,--},G']$ hence $G'$ must be the derivative of $\pi_{13}^*(W)$ by
$\tilde{X}_{1,--}$ that is $G'=Y_1-\Ad^{-1}_{g_{13}}(Y_3).$ That is  after the eliminations $Y_2,\tilde{X}_{1,--}$ we obtain
the complex
\begin{equation}\label{eq: last step red}
 \bar{\pi}^*\left(\bar{\pi}_{12}^*(\calF)\otimes\bar{\pi}_{23}^*(\calG)\right)\otimes [Y_2 -\Ad_{g_{23}}(Y_3)_{++},0]\otimes \pi_{13}^*(CC)
\end{equation}
as we claimed in the proposition.
%Thus if we apply functor $\Hom_{U\frn^{(2)}}(\CE_{\frn^{(2)}},\bullet)$ and take $T^{(2)}$ invariants we obtain $\Phi(\calF\bar{\star}\calG)$ as expected.

%In our argument we did not track the changes of $B^3$-equivariant structure as we were applying elementary transformations to our matrix factorizations but it
%is elementary to check that the $B^3$ structure on (\ref{eq: last step red}) is the expected one.

%\end{proof}


\def\frp{\mathfrak{p}}
\section{Induction functors}\label{sec: inc}

The main goal of this section is to construct induction functors:
$$ \SymbolPrint{\ind_k}: \MF_{B_k^2}(\calX_2^\circ(G_k),W)\times \MF_{B_{n-k}^2}(\calX^\circ_2(G_{n-k}),W)\to \MF_{B_n^2}(\calX_2^\circ(G_n),W),$$
$$ \SymbolPrint{\indb_k}: \MF_{B_k^2}(\calXr_2(G_k),\Wr)\times \MF_{B_{n-k}^2}(\calXr_2(G_{n-k}),\Wr)\to \MF_{B_n^2}(\calXr_2 (G_n),\Wr).$$
 and to show that these functors are intertwined by the Knorrer functors.
\begin{proposition}\label{prop: ind phi} We have
$$ \ind_k\circ (\Phi_k\times \Phi_{n-k} ) =\Phi_n\circ \ind_k.$$
\end{proposition}
We also show that the induction functor is the convolution algebra homomorphism.
\begin{proposition}\label{prop: ind homo} For any
  \[\calF_1,\calG_1\in \MF_{B_k^2}(\calX^\circ_2(G_k),W), \quad \calF_2,\calG_2\in \MF_{B_{n-k}^2}(\calX^\circ_2(G_{n-k}),W).\]
  We have
$$ \ind_k(\calF_1,\calF_2)\star\ind_{k}(\calG_1,\calG_2)=\ind_k(\calF_1\star\calG_1,\calG_1\star\calG_2).$$
\end{proposition}

Finally, we define the generators of the braid group on two strands and using the induction functor explain how we construct the generators for the braid group on
arbitrary number of strands.

\subsection{Construction of the induction functor}
Let us set notations for the parabolic subgroups of $G_n$. We denote by $P_k\subset G_n$  the subgroup with Lie algebra generated by
$E_{ij}$, $i\le j$ and $E_{i+1,i}$, $i\ne k$. We denote the corresponding Lie algebra as $\SymbolPrint{\frp_k}:=\Lie(P_k)$.
There is a natural homomorphism $\SymbolPrint{i_k}: P_k\to G_n$.



We define $\calX^\circ_2(P_k)\subset\calX^\circ_2(G_n)$ as a locus of $(X,g_{12},Y_1,Y_2)$ such that $g_{12}\in P_k$ and $X\in \frp_k$.
Let $i_k$ be the natural  map $\calX^\circ_2(P_k)\to \calX_2^\circ(G_n)$. The image of $i_k$ is defined by the equations
$\SymbolPrint{\Pi_{--}^k}(g_{12})=0$, $\Pi_{--}^k(X)=0$  where $\Pi^k_{--}$ is the projection from the space of $n\times n$ matrices to its $ij$ entries $i\in [k+1,n]$, $j\in [1,k]$.

The ideal generated by $\Pi_{--}^k(g_{12})$, $\Pi_{--}^k(X)$  is $B_n^2$-invariant and these elements form a regular sequence. Hence there is a well-defined functor
$$ i_{k*}:\MF_{B^2}(\calX^\circ_2(P_k),W)\to \MF_{B^2}(\calX_2^\circ(G_n),W)$$
induced by the embedding.

Let us denote by $\SymbolPrint{\Pi_{++}^k}:\frg\to \frg$ the linear projection from the space of $n\times n$ matrices to its  $ij$ entries $i\in [1,k]$, $j\in [k+1,n]$. In particular
the image of $(\mathrm{Id}-\Pi_{++}^k)(\frp_k)$ is naturally identified with $\frg_k\oplus\frg_{n-k}$, we denote this map by $p_k:\frp_k\to \frg_k\oplus \frg_{n-k}$.
Similarly, we have a map $\SymbolPrint{p_k}:P_k\to G_k\times G_{n-k}$.

By applying $p_k$ to all factors of the space $\calX_2^\circ(P_k)$ we obtain the map
$p_k:\calX^\circ_2(P_k)\to \calX^\circ_2(G_k)\times\calX^\circ_2(G_{n-k})$.
%Let $B_{k,n-k}\subset B_n$ be the image of the map $p_k$ which  is a homomorphism of groups. The space $\calX_2^\circ(G_{k,n-k})$ has a natural $B_{k,n-k}^2$ action
%that makes the map $p_k$ $(B_{k,n}^2)$-equivariant. Moreover, 
Using homomorphism $p_k\colon B_n\to B_{k,n-k}$ we can extend $\frn_k^2\times \frn_{n-k}^2$-equivariant structure of $\calX_2(G_k)\times \calX_2(G_{n-k})$ to the
$\frn_n^2$ equivariant structure and the map $p_k$ is $\frn_n^2$-equivariant in this setting.
Moreover, since $T_n=T_k\times T_{n-k}\times \mu$ where $\mu$ is an finite subgroup of the group of roots of unity, we have the following functor:
$$ p_k^*: \MF_{B_k^2}(\calX^\circ_2(G_k),W)\times \MF_{B_{n-k}^2}(\calX^\circ_2(G_{n-k},)W)\to \MF_{B_n^2}(\calX^\circ_2(P_n),W).$$

%Similary to the above case we construct $B_n^2$ action on $\calX_2^\circ(G_n)\times\calX_2^\circ(G_{n-k})$.
%The map $\det$ induces inclusion map  $\calX_2^\circ(G_n)\times\calX_2^\circ(G_{n-k})\to \calX^\circ_2(G_{k,n-k})$ and this map is $B_n^2$-equivariant inclusion.
%Thus
%we have well defined push forward $\det_*: \MF_{B_k^2}(\calX_2^\circ(G_k),W)\times \MF_{B_{n-k}^2}(\calX_2^\circ(G_{n-k})\to \MF_{B_n^2}(\calX^\circ_2(G_{k,n-k}),W)$.

The induction functor is defined as composition of maps:
$$  \ind_k:=i_{k*}\circ p_k^*$$


\subsection{Computational aspects of push-forward}\label{ssec: comp ins}
In our general construction of the push forward $i_{k*}(\calF)$, $\calF=(M,D,d_l,d_r)$  we make a choice of the extension $\tilde{D}$ of the differential $D$
from sub-variety $\calX^\circ_2(P_k)$ to $\calX^\circ_2(G_n)$. However in our case we have a canonical choice of
the extension because we could use the projection map $\Lie(G_n)\to \Lie(P_{k})$ to construct such an extension and use the remark~\ref{rem: section}, the details are provided below.
Morally, we would like to interpret $i_{k}$ as a section of some map and it is possible if we enlarge our space slightly.


In this subsection we use notation $\adj(A)$ for the adjoint of matrix $A$.
Let us introduce space $\SymbolPrint{\calX^\circ_2(G_n)^\flat}:=\frg_n\times \Lie(G_n) \times \frn_n^{2}$. There is
a natural embedding $$j_n: \calX_2^\circ(G_n)\to \calX_2^\circ(G_n)^\flat$$ induced by the embedding $G_n\to \mathfrak{gl}_n$. If we define
\[W^\flat(X,g,Y_1,Y_2)=
\Tr(X( Y_1\det- gY_2\adj(g))),\quad \det=\det(g)\] then we have
$j^*_n(W\cdot\det)=W^\flat$ and the map \[j^*_n: \MF_{B^2_n}(\calX_2^\circ(G_n)^\flat,W^\flat)\to \MF_{B^2_n}(\calX_2^\circ(G_n)^\flat,W'),\]
where \(W'=W\cdot\det\), is well defined.

Similarly, we define  $\calX^\circ_2(P_k)^\flat$ and $\calXr^\circ_2(P_k)^\flat$.
  The natural  map $$\SymbolPrint{j^k_n}: \calX^\circ_2(P_k)\to \calX^\circ_2(P_k)^\flat$$
induces the pull back map $$j^{k*}_{n}: \MF_{B^2}(\calX^\circ_2(P_k)^\flat,W^\flat)\to \MF_{B^2}(\calX^\circ_2(P_k),W').$$


Since \(W'\) and \(W\) differ by the non-vanishing factor \(\det\), there is an embedding of categories \(\det: \MF_{B^2}(\calX^{\circ}_2,W)\rightarrow\MF_{B^2}(\calX^{\circ}_2,W')\):
\[\det: [M_0\xrightarrow{D_0} M_1\xrightarrow{D_1} M_1]\mapsto [M_0\xrightarrow{D_0\cdot \det} M_1\xrightarrow{D_1}M_0]. \]
We the inverse functor we denote by \(\det^{-1}\) and we use the same notation for the functors between the categories related to the parabolic versions of
the categories.


There is no natural map between $\calX^\circ_2(G_n)$ and $\calX^\circ_2(P_k)$ but there is a natural linear projection $\SymbolPrint{\rho_k}: \calX^\circ_2(G_n)^\flat\to
\calX_2^\circ(P_k)^\flat$ and
there is a pull back map $$\rho_k^*: \MF(\calX_2^\circ(P_k)^\flat,W^\flat)\to \MF(\calX_2^\circ(G_n)^\flat,\rho_k^*(W^\flat)).$$

There is also a natural inclusion map $i_k: \calX^\circ_2(P_k)^\flat\to \calX^\circ_2(G_n)^\flat$ which is a section of the projection
$\rho_k$: $\rho_k\circ i_{k}=id$.
%
Hence there is $\calF'\in \MF_{B^2_n}(\calX^\circ_2(P_k)^\flat,W^\flat)$ such that $j^{k*}_n(\calF')=\det(\calF)$
and we define \[\tilde{\calF}=(\tilde{M},\tilde{D},\tilde{\partial}_l,\tilde{\partial}):=\mathrm{det}^{-1}\circ j^{*}_n\circ \rho_k^*(\calF')\in \MF(\calX^\circ_2(G_n),j_n^{*}\circ{\rho}_{k}^*(W^\flat)\cdot\mathrm{det}^{-1}).\]
%From now on we abbreviate $j_n^*\circ{\rho}_{k\dots m}^*(\Wr^\flat)$ as $\Wr_{k\dots m}$.
By construction we have $\tilde{\calF}|_{\calX^\circ_2(P_k)}=\calF$. Thus the matrix factorization  $\tilde{\calF}$ provides the first step of the construction of the push forward
$i_{k*}(\calF)$ from section~\ref{ssec: eq push f}. By completing the second step of the construction from the section~\ref{ssec: eq push f} we obtain the matrix factorization:
%so this complex satisfies the assumption of the extension lemma~\ref{lem: ext MF+eq} and the push-forward matrix factorization $i_{k*}(\calF)\in \MF_{B^2}(\calX^\circ_2(G_n),W)$.
% such that
$$ i_{k*}(\calF)=(\tilde{M},d_K+\tilde{D}+d^-,\partial_l,\partial_r),$$
where $d_K: \tilde{M}\otimes \Lambda^{\bullet} V\to \tilde{M}\otimes \Lambda^{\bullet-1} V$ is the Koszul differential for the
defining ideal of $\calX_2^\circ(P_k)\subset \calX_2^\circ(G_n)$ and $d^-=\sum_{i<j} d^-_{ij}$,
$d^-_{ij}:\tilde{M}\otimes\Lambda^i V\to \tilde{M}\otimes \Lambda^{j} V$ is the complementary differential: $(d_k+\tilde{D}+d^-)^2=\Wr$

Note that because of our choice of extension $\tilde{\calF}$ the differentials $d^-_{ij}$ vanish unless $j=i+1$, we give more details below.
Define a linear operator operator $\Pi^k_{+}:\frg_n\to \mathfrak{p}_k$ which is a natural projection. Respectively, we have $\Pi_{--}^k(A)=A-\Pi_+^{k}(A)$.
The difference $\delta_kW:=W-j_n^*\circ\rho_n^{k*}(W^\flat)\cdot \det^{-1}$ could be written as sum of two terms:
$$\delta_kW_1:=\Tr(\Pi_{--}^k(X)\cdot \Pi_{++}^k(Y_1-\Ad_{g_{12}}(Y_2))),$$
$$\delta_kW_2:=\Tr(\Pi_+^k(X)\cdot
\Pi^k_{-}(\Ad_{g_{12}}(Y_2)-\rho_k(g_{12})(Y_2)\adj(\rho_k(g_{12}))\det(g_{12})^{-1}  ),$$
where we use convention $\mathrm{Id}=\Pi^k_{-}+\Pi^k_{++}$.

Since $\delta_k W_1\in (\Pi^k_{--}(X))$ and $\delta_kW_2\in (\Pi^k_{--}(g_{12}))$ we have
$$ i_{k*}(\calF)^\sharp=\tilde{\calF}\otimes [\Pi_{--}^k(X), \Pi_{++}^k(Y_1-\Ad_{g_{12}}(Y_2))]\otimes  \mathrm{K}^{\delta_k W_2} (\Pi_{--}^{k}(g_{12})).$$
That is the differentials \(d_{ij}^-\)   are the negative differentials of the Koszul matrix factorization and hence have above mentioned property.



\subsection{Induction functor in reduced case}

%Lets fix $k<m$ and define $P_{k\dots m}\subset G_n$ to be the subgroup defined by the conditions that $ij$-entry of the matrix vanishes if $j>i$ and $i,j\notin [k,m]$ and $\mathbf{P}_{k\dots m}\subset \bfG_n$  is the
%subgroup defined by the same equations. Let us use notation $\mathfrak{p}_{k\dots m}$ for $\mathrm{Lie}(\mathbf{P}_{k\dots m})$.
Let us
define $\calXr_2(P_k):=\frb\times P_k\times \frn$ and $\calXr_2(P_k)^\flat:=\frb\times\mathrm{Lie}(P_k)\times\frn$.
We also use notation $\Wr$ for the restriction of the potential $\Wr$ on the subspace $\calXr_2(P_{k})$.

There is a natural embedding $\bar{i}_k: \calXr_2(P_k)\to \calXr_2(G_n)$.
The map of the rings $\bar{i}_{k}: \CC[\calXr_2(G_n)]\to \CC[\calXr_2(P_k)]$ is the projection map with the kernel $(\Pi^k_{--}(g))$ and it is easy to see that these generators
form a regular sequence. Thus we have a well-defined functor:
$$ \bar{i}_{k*}: \MF_{B^2}(\calXr_2(P_k),\Wr)\to \MF_{B^2}(\calXr_2(G_n),\Wr).$$
We also have the projection map $\bar{p}_k: \calXr_2(P_k)\to \calXr_2(G_k)\times\calXr_2(G_{n-k})$.


The map $\bar{p}_k$ is $B_{k}^2\times B_{n-k}^2$-equivariant and we extend $B_{k}^2\times B_{n-k}^2$-equivariant structure to
$B_n^2$-equivariant structure to obtain  an induced map $$\bar{p}^*_k:
\MF_{B_k^2}(\calXr_2(G_k),\Wr)\times\MF_{B_{n-k}^2}(\calXr_2(G_{n-k}),\Wr) \to \MF_{B_n^2}(\calXr_2(P_k),\Wr).$$  Thus we define
$$\indb_k:=\bar{i}_{k*}\circ \bar{p}_k^*.$$


Now let us discuss an analogs of the map $\rho_k$ for reduced spaces and the corresponding properties of the push forward $\bar{i}_{k*}$.
%\subsection{Computational aspects of push-forward}
%In our construction of the push forward $i_{k\dots m*}(\calF)$, $\calF=(M,D,d_l,d_r)$  we make a choice of the extension $\tilde{D}$ of the differential $D$
%from sub-variety $\calXr_2(P_{k\dots m})$ to $\calXr_2(G_n)$. However in our case we have somewhat canonical choice of
%the extension because we could use the projection map $\Lie(\bfG_n)\to \Lie(\mathbf{P}_{k\dots m})$ to construct such an extension and use the remark~\ref{rem section}, the details are provided below.
%Morally, we would like to interpret $\bar{i}_{k\dots m}$ as a section of some map and it is possible if we enlarge our space slightly.
%In this subsection we use notation $\adj(A)$ for the adjoint of matrix $A$.
Consider a space $\calXr_2^{\flat}(G_n):=\frb\times \mathfrak{gl}_n\times \frn$. There is
a natural embedding $\bar{j}_n: \calXr_2(G_n)\to \calXr_2(G_n)^\flat$ induced by the embedding $G_n\to \mathfrak{g}_n$. If we define $\Wr^\flat(X,g,Y)=\Tr(X( gY\adj(g)))$ then we have
$\bar{j}^*_n(\Wr')=\Wr^\flat$, \(\Wr'=\Wr\cdot\det(g)\) and map
$$\bar{j}^*_n: \MF_{B^2_n}(\calXr_2(G_n)^\flat,\Wr^\flat)\to
\MF_{B^2_n}(\calXr_2(G_n),\Wr')$$ is well-defined.

Similarly we define  $\calXr_2(P_{k})$ and $\calXr_2(P_{k})^\flat$.
  The natural open embedding $\bar{j}^k_n: \calXr_2(P_{k})\to \calXr_2(P_{k})^\flat$
  induces the pull back map $$\bar{j}^{k*}_n: \MF_{B^2}(\calXr_2(P_k),\Wr')\to \MF_{B^2}(\calXr_2(P_{k})^\flat,\Wr^\flat).$$
  
 There is a natural linear projection $$\bar{\rho}_k: \calXr_2(G_n)^\flat\to \calXr_2(
P_{k})^\flat$$ and
 a pull back map $$\bar{\rho}_k^*: \MF(\calXr_2(P_{k})^\flat,\Wr^\flat)\to \MF(\calXr(G_n)^\flat,\rho_{k}^*(\Wr^\flat)).$$

%There also a natural inclusion map $j_{k\dots m}: \calXr_2(\mathbf{P}_{k\dots m})^\flat\to \calXr_2(\bfG_n)^\flat$ which is a section of the projection
%$\rho_{k\dots m}$: $\rho_{k\dots m}\circ j_{k\dots m}=id$.

 We can choose $\calF'\in \MF_{B^2}(\calX_2(G_n)^\flat,\Wr)$ such that $\bar{j}^{k*}_n(\calF')=\det(\calF)$
 and we define
 \[\tilde{\calF}:=\mathrm{det}^{-1}\circ\bar{j}^{k*}_n\circ \bar{\rho}_{k}^*(\calF')\in \MF(\calX_2(G_n),\bar{j}_n^{k*}\circ\bar{\rho}_{k}^*(\Wr^\flat)\cdot\mathrm{det}^{-1}).\]
Let $\delta_k\Wr=\Wr-\bar{j}_n^{k*}\circ\bar{\rho}_k^*(\Wr^\flat)/\det$, then analogously to the non-reduced case we have:
%From now on we abbreviate $j_n^*\circ{\rho}_{k\dots m}^*(\Wr^\flat)$ as $\Wr_{k\dots m}$.
%By construction we have $\tilde{\calF}|_{\calXr_2(P_{k\dots m})}=\calF$ thus
%this complex satisfies the assumption of the extension lemma~\ref{lem: ext MF+eq} and the push-forward matrix factorization $i_{k\dots m*}(\calF)\in \MF_{B^2}(\calXr_2(P_{k\dots m},\Wr)$.
% such that
%$$ i_{k\dots m*}(\calF)=(\tilde{M},d_K+\tilde{D}+d^-,\partial_l,\partial_r),$$
%where $d_K: \tilde{M}\otimes \Lambda^{\bullet} V\to \tilde{M}\otimes \Lambda^{\bullet-1} V$ is the Koszul differential for the ideal $J_{k\dots m}$ and $d^-=\sum_{i<j} d^-_{ij}$,
%$d^-_{ij}:\tilde{M}\otimes\Lambda^i V\to \tilde{M}\otimes \Lambda^{j} V$ is the complementary differential: $(d_k+\tilde{D}+d^-)^2=\Wr$
%
%Let us also notice that because of our choice of extension $\tilde{F}$ the differentials $d_{ij}$ vanish unless $j=i+1$, let us give more details below.
%Let us define a linear operator operator $\Pi^{k\dots m}_+:\mathfrak{gl}(n)\to \mathfrak{p}_{k\dots m}$ which is a natural projection. Respectively, we can define $\Pi_-^{k\dots m}(A)=A-\Pi_+^{k\dots m}(A)$.
%Since $\Wr-\Wr_{k\dots m}\in (\Pi_-^{k \dots m}(g_{12}))$ there is a well define Koszul complex $\mathrm{K}^{\Wr-\Wr_{k\dots m}}(\Pi_-^{k\dots m}(g_{12}))\in \MF(\calXr_2,\Wr-\Wr_{k\dots m})$ and
$$ \bar{i}_{k*}(\calF)^\sharp=\tilde{\calF}\otimes \mathrm{K}^{\delta_k\Wr} (\Pi_{--}^{k}(g_{12})).$$


\subsection{Proof of proposition~\ref{prop: ind phi}}
Note that the statement of the proposition refers to the Knorrer functor
$$\Phi_n:\quad \MF_{B_n^2}(\calXr_2(G_n),\Wr)\to \MF_{B_n^2}(\calX^\circ_2(G_n),W),$$
which is constructed the in same way as our usual functor, the only adjustment being
%with only adjustment that consists of
the replacement of the
space $\widetilde{\calX}_2$ by the space $\widetilde{\calX}_2^\circ$ and the appropriate adjustment of the maps.

The statement of the proposition is equivalent to the isomorphism between two functors relating the following categories:
$$\MF_{B_k^2}(\calXr(G_k),\Wr)\times \MF_{B_{n-k}^2}(\calXr_2(G_{n-k}),\Wr)\longrightarrow \MF_{B_n^2}(\calX^\circ_2(G_n),W).$$
The first functor is $ i_{k*}\circ p_k^*\circ (j^x\times j^x)_*\circ (\pi_y\times\pi_y)^*$ and the second is $j^x_*\circ \pi_y^*\circ\bar{i}_{k*}\circ \bar{p}_k^*$.

We begin with the first functor. Using maps $p_k:\calX^\circ_2(P_k)\to \calX^\circ_2(G_k)\times \calX^\circ_2(G_{n-k})$ and
$j^x\times j^x: \widetilde{\calX}_2^\circ(G_k)\times \widetilde{\calX}_2^\circ(G_{n-k})\to \calX^\circ_2(G_k)\times \calX^\circ_2(G_{n-k})$ we can form
 a product of $\calX^\circ_2(P_k)$ and $ \widetilde{\calX}_2^\circ(G_k)\times \widetilde{\calX}_2^\circ(G_{n-k})$ over $\calX^\circ_2(G_k)\times \calX^\circ_2(G_{n-k})$ which is
 equal $\frb_n\times P_k\times \frn_n^2$. Thus we obtain a commuting diagram
 $$ \begin{tikzcd}
 \frb_n\times P_k\times \frn_n^2 
 %\arrow[r, dotted,"j^x"]
 \arrow[dotted]{r}{j^x} 
 \arrow[dotted]{d}{p_k}   & \calX^\circ_2(P_k)\arrow{d}{p_k}\\
\widetilde{\calX}_2^\circ(G_k)\times \widetilde{\calX}_2^\circ(G_{n-k})\arrow{r}{j^x\times j^x} &                                  \calX^\circ_2(G_k)\times \calX^\circ_2(G_{n-k})
 \end{tikzcd}$$
%$$ \begin{tikzcd}
% \frb_n\times P_k\times \frn_n^2 \arrow[r, dotted,"j^x"]\arrow[d,dotted,"p_k"]   & \calX^\circ_2(P_k)\arrow[d,"p_k"]\\
%\widetilde{\calX}_2^\circ(G_k)\times \widetilde{\calX}_2^\circ(G_{n-k})\arrow[r,"j^x\times j^x"] &                                  \calX^\circ_2(G_k)\times \calX^\circ_2(G_{n-k})
% \end{tikzcd}$$
where the horizontal dotted arrow map is a natural embedding $j^x: \frb_n\times P_k\times \frn_n^2\to \frg_n\times P_k\times \frn_n^2$ while the vertical one is
$$p_k(X,g,\vec{Y})=(p_k(X)',p_k(g)',p_k(Y_1)',p_k(Y_2)')\times(p_k(X)'',p_k(g)'',p_k(Y_1)'',p_k(Y_2)''),$$
where \(\vec{Y}=(Y_1,Y_2).\)

The maps in the diagram are either projections or regular embeddings thus we have  the base change formula:
$ p_k^*\circ (j^x\times j^x)_*=j^x_*\circ p_k^*.$
Thus we have $$ i_{k*}\circ p_k^*\circ (j^x\times j^x)_*\circ (\pi_y\times\pi_y)^*=(i_k\circ j^x)_*\circ (\pi_y\times \pi_y\circ p_k)^*.$$


Analysis of the second functor is similar. As in the previous case we have the the commuting diagram:

$$ \begin{tikzcd}
 \frb_n\times P_k\times \frn_n^2 \arrow[dotted]{r}{i_k}\arrow[dotted]{d}{\pi_y}   & \widetilde{\calX}^\circ_2(G_n)\arrow{d}{\pi_y}\\
\calXr_2^\circ(P_k)\arrow{r}{\bar{i}_k} &                                  \calXr_2(G_n)
 \end{tikzcd}
 $$
where the dotted arrow maps are: $i_k: b_n\times P_k\times \frn_n^2\to \widetilde{\calX}_2^\circ(G_n)$ is induced my $i_k:P_k\to G_n$;
$\pi_y(X,g,Y_1,Y_2)=(X,g,Y_2)$.
%
In this case  we have base change equation: $\pi_y^*\circ \bar{i}_{k*}=i_{k*}\circ \pi_y^*$ and
$$j^x_*\circ \pi_y^*\circ\bar{i}_{k*}\circ \bar{p}_k^*=(j^x\circ i_k)_*\circ (\bar{p}_k\circ \pi_y)^*.$$

Finally note that the map $i_k\circ j^x$ is equal the map to $j^x\circ i_k$ and the map $\pi_y\times \pi_y\circ p_k$ is equal to the map
$\bar{p}_k\circ \pi_y$.

\subsection{Proof of proposition~\ref{prop: ind homo}}
We have to show that
$$ i_{k*}\circ p_k^*(\calF)\star i_{k*}\circ p_{k}^*(\calG)= i_{k*}\circ p_{k}^*(\calF\star\calG),$$
where $\calF=\calF_1\boxtimes\calF_2$, $\calG=\calG_1\boxtimes\calG_2\in \MF_{B_k^2}(\calX^\circ_2(G_k),W)\times\MF_{B_{n-k}^2}(\calX_2^\circ(G_{n-k}),W).$
The LHS of the later formula is the result of the application of the functor $\CE_{\frn^{(2)}_n}$ combined with extraction of  $T^{(2)}$ invariant part to the matrix factorization
$C$ such that:
%
\begin{multline}\label{eq: long mf prod}
C^\sharp=\tilde{\calF}\otimes [\Pi_{--}^k(X), \Pi_{++}^k(Y_1-\Ad_{g_{12}}(Y_2))]\otimes  \mathrm{K}^{\pi_{12}^*(\delta_k W_2)} (\Pi_{--}^{k}(g_{12}))\otimes\\
\tilde{\calG}\otimes [\Pi_{--}^k(\Ad_{g_{12}}^{-1}(X)), \Pi_{++}^k(Y_2-\Ad_{g_{23}}(Y_3))]\otimes  \mathrm{K}^{\pi_{23}^{\circ*}(\delta_k W_2)} (\Pi_{--}^{k}(g_{23})).
\end{multline}
%
%$$C^\sharp= \mathrm{K}^{\bar{\pi}_{12}^*(\Wr-\Wr_{k\dots m})}(\Pi_-^{k\dots m}(g_{12}))\otimes  \mathrm{K}^{\bar{\pi}_{23}^*(\Wr-\Wr_{k\dots m})}(\Pi_-^{k\dots m}(g_{23}))
%\otimes \bar{\pi}_{12}^*(\tilde{\calF})\otimes \bar{\pi}_{23}^*(\tilde{\calG}),
%$$
with  $\tilde{\calF}= j_n^{k*}\circ \rho_{k}^*(\calF')$ and $\tilde{\calG}=j_n^{k*}\circ\rho_{k}^*(\calG')$ in conventions of subsection~\ref{ssec: comp ins}:
$j_n^{k*}(\calF')=p^*_{k}(\calF)$, $j_{n}^{k*}(\calG')=p^*_{k}(\calG)$.

Let us fix coordinates $X,g_{23},g_{13},Y_1,Y_2,Y_3$ on $\calX^\circ_3(G_n)$.
We use  the factor $[\Pi_{--}^k(X), \Pi_{++}^k(Y_1-\Ad_{g_{12}}(Y_2))]$ to exclude the variables $\Pi_{--}^k(X)$ from the other factors.
Since $\Pi_{--}^k(X)$  is $B_n^2$-equivariant, the exclusion process uses only the row transformations of the first
kind and we do not need to worry about the admissibility conditions.
The elimination process changes the factor in  (\ref{eq: long mf prod}) in the following manner:
$$\tilde{\calF}\to\tilde{\calF}_1:= \tilde{\calF}|_{\Pi_{--}^k(X)=0},\quad
     \tilde{\calG}\to\tilde{\calF}_1:= \tilde{\calF}|_{\Pi_{--}^k(\Ad_{g_{12}}^{-1}(X))=0}$$
$$  [\Pi_{--}^k(X), \Pi_{++}^k(Y_1-\Ad_{g_{12}}(Y_2))]\to  [\Pi_{--}^k(X), R_1],$$
$$ [\Pi_{--}^k(\Ad_{g_{12}}^{-1}(X)), \Pi_{++}^k(Y_2-\Ad_{g_{23}}(Y_3))]\to [0,\Pi_{++}^k(Y_2-\Ad_{g_{23}}(Y_3))],$$
$$ \mathrm{K}^{\pi_{12}^{\circ*}(\delta_k W_2)} (\Pi_{--}^{k}(g_{12}))
\to  \mathrm{K}^{\pi_{12}^*(\delta_k W_2)} (\Pi_{--}^{k}(g_{12}))|_{\Pi_{--}^k(X)=0},$$
$$  \mathrm{K}^{\pi_{12}^{\circ*}(\delta_k W_2)} (\Pi_{--}^{k}(g_{23}))\to  \mathrm{K}^{\pi_{23}^{\circ*}(\delta_k W_2)} (\Pi_{--}^{k}(g_{23}))|_{\Pi_{--}^k(X)=0}$$
We denote the resulting equivariant matrix factorization as $C_1$. The non-equivariant limit $C_1^\sharp$ is the product of the above listed terms. Moreover, since the only term in
$C_1^\sharp$ that depends on $\Pi_{--}^k(X)$ is the Koszul matrix factorization $[\Pi_{--}^k(X),R_1]$, we conclude that $R_1$ is the derivative of the potential $\pi_{13}^{\circ*}(W)$ by
variables $\Pi_{--}^k(X)$ which equals $\Pi_{++}^k(Y_1-\Ad_{g_{13}}(Y_3))$.

Next let us observe that using the row transformations of the first kind we can show that the following equivalence of the matrix factorization
 \begin{multline*}
 \mathrm{K}^{\pi_{12}^{\circ,*}(\delta_k W_2)} (\Pi_{--}^{k}(g_{12}))\otimes \mathrm{K}^{\pi_{23}^{\circ*}(\delta_k W_2)} (\Pi_{--}^{k}(g_{23}))|_{\Pi_{--}^k(X)=0}\\
 = \mathrm{K}^{\pi_{13}^{\circ*}(\delta_k W_2)} (\Pi_{--}^{k}(g_{13}))\otimes \mathrm{K}^{0} (\Pi_{--}^{k}(g_{23}))|_{\Pi_{--}^k(X)=0}
 \end{multline*}

 Finally, we use the Koszul complex $\mathrm{K}^{0} (\Pi_{--}^{k}(g_{23}))|_{\Pi_{--}^k(X)=0}$ to exclude the variables $\Pi_{--}^k(g_{23})$ from the product; for that
 we use the first type row transformations so we do not need to worry about the admissibility. The final result of the sequence of row transformations is the
 equivariant matrix factorizations $C_2$ such that its non-equivariant limit $C_2^\sharp$ is the product the following factors:
 $$\tilde{\calF}_2:= \tilde{\calF}|_{\Pi_{--}^k(X)=0,\Pi_{--}^k(g_{23})=0},\quad
     \tilde{\calG}_2:= \tilde{\calF}|_{\Pi_{--}^k(\Ad_{g_{12}}^{-1}(X))=0, \Pi_{--}^k(g_{23})=0}$$
 $$[\Pi^k_{--}(X),Y_1-\Ad_{g_{13}}(Y_3)],\quad     [0,\Pi_{++}^k(Y_2-\Ad_{g_{23}}(Y_3))],$$
$$ \mathrm{K}^{\pi_{13}^{\circ*}(\delta_k W_2)} (\Pi_{--}^{k}(g_{13}))\otimes \mathrm{K}^{0} (\Pi_{--}^{k}(g_{23}))|_{\Pi_{--}^k(X)=0}.$$
Since the potential $\pi_{13}^{\circ*}(\delta_k W_2)$ does not depend on $\Pi_{--}^{k}(X)$ the last product is actually just
 $$ \mathrm{K}^{\pi_{13}^{\circ*}(\delta_k W_2)} (\Pi_{--}^{k}(g_{13}))\otimes \mathrm{K}^{0} (\Pi_{--}^{k}(g_{23})).$$

 Let us introduce auxiliary space $\SymbolPrint{\calZ_{con}(G_n)}:=\frg_n\times P_k\times G_n\times \frn_n\times (\frn_k\times \frn_{n-k})\times \frn_n$
 and $\SymbolPrint{\calZ_{con}(P_k)}:=\frp_n\times P_k\times P_k\times \frn_n\times (\frn_k\times \frn_{n-k})\times \frn_n$. The variety $\calZ_{con}(G_n)$ is a subspace of
 $\calX^\circ_3$ defined by equations $\Pi_{++}^k(Y_2-\Ad_{g_{23}}(Y_3))=0$, $\Pi_{--}^k(g_{23})=0$, thus we can define the map $\pi_{13}^\circ:\calZ_{con}(G_n)\to
 \calX_2^\circ(G_n)$ by the restriction.

   We can contract  differentials in complexes $ [0,\Pi_{++}^k(Y_2-\Ad_{g_{23}}(Y_3))]$,
 $\Pi_{--}^{k}(g_{23})$ to obtain homotopy equivalence of matrix factorizations over $\pi_{13}^{\circ,*}(\CC[\calX^\circ_2(G_n)])$:
 $$ C_2\sim C_3,$$
 where $C_3\in \MF_{B_n^2}(\calZ_{con}(G_n),\pi_{13}^\circ(W))$. Moreover, since we only used admissible row transformations the matrix factorization $C_3$ has a
 structure of quasi-Koszul matrix factorization:
 $$ C_3=(\tilde{M}\otimes \Lambda^\bullet V, d^++D+d^-,\partial+\partial^-),$$
 where $V$ is space dual to span of equations $\Pi_{--}^k(X),\Pi_{--}^k(g_{13})$, $d^+$ is the corresponding Koszul differential
$d^+:\tilde{M}\otimes\Lambda^\bullet V\to \tilde{M}\otimes \Lambda^{\bullet-1} V,$ and $d^-: \tilde{M}\otimes\Lambda^\bullet V\to \tilde{M}\otimes \Lambda^{\bullet+1} V$
the completing differentials and \(\partial^-=\sum_{i>j} \partial_{ij}^-\), \[\partial_{ij}: \Lambda^i (\frn_n^2)\otimes \tilde{M}\otimes \Lambda^\bullet V\to
\Lambda^j (\frn_n^2)\otimes \tilde{M}\otimes \Lambda^{\ge\bullet} V, \quad (\tilde{M},D,\partial)\in \Mod_{per}^{B_n^3}(\calZ(G_n)).\]

The locus inside $\calZ_{con}(G_n)$ defined by the equations $\Pi_{--}^k(X)=0,\Pi_{--}^k(g_{13})=0$ is isomorphic to $\calZ_{con}(P_n)$, by construction
$$(\tilde{M},D,\partial)|_{\Pi_{--}^k(X)=0,\Pi_{--}^k(g_{13})=0}=p_k^*(\pi_{12}^{\circ*}(\calF)\otimes\pi_{23}^{\circ*}(\calG)),$$
where $p_k: \calZ_{con}(P_k)\to \calX^\circ_3(G_k)\times\calX^\circ_3(G_{n-k})$ is the natural projection.
Hence we can use uniqueness lemma~\ref{lem: uniq MF+eq} to imply that we have an isomorphism of $B_n^3$ equivariant matrix factorizations
$$ C_3=i_*(p_k^*(\pi_{12}^{\circ*}(\calF)\otimes\pi_{23}^{\circ*}(\calG)),$$
where $i$ is map $\calZ_{con}(P_k)\to \calZ_{con}(G_n)$.

The Lie algebra $\frn^{(2)}$ does not act on variables $\Pi^k_{--}(X)$ and $\Pi_{--}^k(g_{13})$ hence we have
$$ \CE_{\frn^{(2)}}(C_3)=i_*(\CE_{\frn^{(2)}}(p_k^*(\pi_{12}^{\circ*}(\calF)\otimes \pi_{23}^{\circ*}(\calG)))).$$

Let us denote by $I_k$ the kernel of the projection homomorphism $P_k\to B_k\times B_{n-k}$ and let $I'_k$ be its unipotent part. Hence we have the short exact sequence of
Lie algebras:
$$0\to \Lie{(I'_k)}\to \frn_n \to \frn_k\oplus \frn_{n-k}\to 0$$
and  we can use Hochschild-Serre spectral sequence to compute functor $\CE_{\frn_n^{(2)}}$:
$$ \CE_{\frn_{k}^{(2)}\oplus \frn_{n-k}^{(2)}}(\CE_{\Lie(I'_{k})}(C_3))\Rightarrow \CE_{\frn_n^{(2)}}(C_1). $$


Since the Lie algebra $\frn^{(2)}$ only acts on on the second copy of $P_k$ in the product
$\calZ_{con}(P_n)=\frp_n\times P_k\times P_k\times \frn_n\times (\frn_k\times \frn_{n-k})\times \frn_n$
and the differentials in $p_k^*(\pi_{12}^{\circ*}(\calF)\otimes \pi_{23}^{\circ*}(\calG))$ are also $I'_k$ invariant,
the computation of $\CE_{\Lie(I'_k)}(p_k^*(\pi_{12}^{\circ*}(\calF)\otimes \pi_{23}^{\circ*}(\calG)))$ reduces to the computation
of
$$\textup{H}_{\Lie}^*(\Lie(I'_k),\CC[\calZ_{con}(P_k)])=\CC[\frn_n\times P_k\times \frn_n\times (\frn_k\times \frn_{n-k})\times \frn_n]\otimes
\textup{H}_{\Lie}^*(\Lie(I'_k),\CC[P_k]).$$

The last computation is equivalent to the following

\begin{proposition}
$$\CE_{\Lie(I'_k)}(\CC[P_k])=\CC[G_k\times G_{n-k}],$$
where $G_k\times G_{n-k}$ is the subgroup of $P_k$ defined by the equation $\Pi_{++}^k(g)=0$.
\end{proposition}
\begin{proof}
The group $I_k$ acts on $P_k$ by right multiplication. The Lie algebra $\Lie(I'_k)$ is generated by the elements $E_{ij}$, $i\in [1,n]$, $j\in [k+1,n]$ and direct computation
shows that it is a commutative algebra.

Let us denote by $E^\vee_{kl}$ the generators of the algebra $\CC[P_k]$:
$E^\vee_{kl}(X)=X_{kl}.$ Let us denote by $\partial_{ij}$ the differential operator corresponding to the action $E_{ij}\in \Lie(I_k')$.
The action of these generators of $\Lie(I'_k)$ is given by:
$$ \partial_{ij}=\sum_{k+1\le l \le n} E_{jl}^\vee\frac{\partial}{\partial E_{il}^\vee}.$$

Let us observe that vector of differential operators $\partial_{i\bullet}$ is related to the vector of differential operators $\frac{\partial}{\partial E_{i\bullet}^\vee}$, $\bullet\in [k+1,n]$ by the matrix
$M=(E^\vee_{lm})_{l,m\in [k+1,n]}$ and this matrix is invertible. Thus the condition $\partial_{ij}(f)=0$ for all $i,j$ is equivalent to the condition
$\frac{\partial f}{\partial E_{ij}^\vee}=0$ for all $i,j\in [k+1,n]$.

Since the operator $\frac{\partial}{\partial E_{ij}^\vee}$ is surjective on $\CC[P_k]$ the statement of the proposition follows.
\end{proof}

The proposition implies that $\CE_{\Lie(I'_k)}(\CC[\calZ_{con}(P_k)])=\CC[\calZ_{con}(G_k\times G_n)]$ where $\calZ_{con}(G_k\times G_n):=\frp_n\times P_k\times G_k\times G_{n-k}\times \frn_n\times (\frn_k\times \frn_{n-k})\times \frn_n$.
Thus the proposition implies that
$$ \CE_{\frn^{(2)}}(C_3)=i'_*(\CE_{\frn_k\oplus \frn_{n-k}}((p'_k)^*(\pi_{12}^{\circ*}(\calF)\otimes\pi_{23}^{\circ*}(\calG)))),$$
where $i': \calZ_{con}(G_k\times G_{n-k})\to \calZ_{con}(G_n)$ and $p'_k:\calZ_{con}(G_k\times G_{n-k})\to \calX^{\circ}_3(G_k)\times \calX^{\circ}_3(G_{n-k}).$

%{\bf We are not careful here with these annoying roots of unity :( }



% The torus $T^{(2)}$ splits (up-to some small group of unity) as a product $T^{(2)}=\CC^*\times T_k^{(2)}\times T_{n-k}^{(2)}$ where
% $\CC^*\subset I_k$ and $G_{k,n}^{\CC^*}=G_k\times G_{n-k}$ and $T_k^{(2)}$ and $T_{n-k}^{(2)}$ are the maximal tori of
% $G_k$ and $G_{n-k}$ respectively.


Hence after extracting the torus invariant part of the last expression we obtain the formula from the statement of the proposition:
 $$ \pi_{13*}^{\circ}(\CE_{\frn^{(2)}}(C_3))^{T^{(2)}}=i^*_k\circ p_k^*( \calF\bar{\star}\calG).$$





\subsection{Transitivity of induction}
In this section we prove the following
\begin{proposition}\label{prop: ind trans} The following functors:
$$\MFs_k\times \MFs_{m-k}\times \MFs_{n-m}\to \MFs_{n},$$
$\MFs_r:=\MF_{G_r\times B_r^2}(\calX_2(G_r),W)$
 are isomorphic
$$ \ind_k\circ \mathrm{Id}\times \ind_{m-k}=\ind_m\circ \ind_k\times \mathrm{Id}.$$
\end{proposition}

Before we start the proof let us fix some notations. We denote by $\SymbolPrint{P_{k,m}}$ a subgroup of $G_n$ with the Lie algebra generated by
$E_{ij}$, $i\le j$ and $E_{i,i+1}$, $i\ne k,m$. Given an element $g\in P_{k,m}$ we denote by $g'$ its $k\times k$ block, by $g''$ its $(m-k)\times (m-k)$ block and
by $g'''$ its $(n-m)\times (n-m)$ block.
 We also use $\SymbolPrint{\frp_{k,m}}$ for $\Lie(P_{k,m})$

 
 We have natural homomorphisms:
$\SymbolPrint{i_{k,m}}: P_{k,m}\to G_n$ and $p_{k,m}: P_{k,m}\to G_k\times G_{m-k}\times G_{n-m}$.
Respectively, we define $\calX_2^\circ(P_{k,m}):=\frp_{k,m}\times P_{k,m}\times \frn_n^2$ and using previously defined maps we construct
maps $\SymbolPrint{p_{k,m}}: \calX_2^\circ(P_{k,m})\to \calX_2^\circ(G_k)\times \calX_2^\circ(G_{m-k})\times \calX_2^\circ(G_{n-m})$ and
$i_{k,m}:  \calX_2^\circ(P_{k,m})\to \calX_2^\circ(G_n)$.
Using these maps we can define the functor
 $$\SymbolPrint{\ind_{k,m}}: \MFs^\circ_{k}\times \MFs^\circ_{m-k}\times \MFs^\circ_{n-m}\to \MFs^\circ_{n},$$
$$
\ind_{k,m}:=i_{k,m,*}\circ p_{k,m}^*,$$
$\MFs^\circ_r:=\MF_{B^2_r}(\calX^\circ_2(G_r),W).$

Also there are also natural maps $p_k: P_{k,m}\to G_k\times P_{m-k}$, $p_m: P_{k,m}\to P_k\times G_{n-m}$ such that $p_{m-k}\circ p_k=p_{k,m}$ and $p_k\circ p_m=p_{k,m}$.

\begin{proof}[Proof of proposition~\ref{prop: ind trans}]
We  will show that
 $$\ind_k\circ \mathrm{Id}\times \ind_{m-k}=\ind_{k,m}=\ind_m\circ \ind_k\times \mathrm{Id}$$
Let us indicate how we imply the first equality, the second equality has a similar proof. Let us notice that LHS is the composition of several functors:
$i_{k*}\circ p_k^*\circ(\mathrm{Id}\times i_{m-k})_*\circ (\mathrm{Id}\times p_{m-k})^*$.  We have  the
following commuting diagram:
$$ \begin{tikzcd}
 \frp_{k,m}\times P_{k,m}\times \frn_n^2 \arrow[dotted]{r}{i_{m-k}}\arrow[dotted]{d}{p_k}   & \calX^\circ_2(P_k)\arrow{d}{p_k}\\
\calX_2^\circ(G_k)\times \calX_2^\circ(P_{m-k})\arrow{r}{1\times i_{m-k}} &                                  \calX^\circ_2(G_k)\times \calX^\circ_2(G_{n-k})
 \end{tikzcd}$$
where the dotted arrow maps are: $p_k(X,g,Y_1,Y_2)=(p_k(X),p_k(g),Y_1,Y_2)$ where $p_k:P_{k,m}\to G_k\times P_{m-k}$ is the natural homomorphism and
$i_{m-k}=i_{m-k}\times i_{m-k}\times \mathrm{Id} \times \mathrm{Id}$.


Finally let us observe that $p_{m-k}\circ p_k= p_{k,m}$ and $i_k\circ i_{m-k}=i_{k,m}$. Now we use the base change to show:
\begin{multline*}
\ind_k\circ \mathrm{Id}\times \ind_{m-k}= i_{k*}\circ p_k^*\circ (\mathrm{Id}\times i_{m-k})_*\circ (\mathrm{Id}\times p_{m-k})^*\\= i_{k*}\circ i_{m-k*}\circ p_k^*\circ p_{m-k}^*=
(i_{k}\circ i_{m-k})_*\circ (p_{m-k}\circ p_k)^*
=\ind_{k,m}.\end{multline*}
\end{proof}

Since the functor $\Phi_n$ provides an embedding of categories we get an immediate corollary:

\begin{corollary} The following functors from 
\[\MF_{B_k^2}(\calXr_2(G_k),\Wr)\times \MF_{B_{m-k}^2}(\calXr_2(G_{m-k}),\Wr)\times \MF_{B_{n-m}^2}(\calXr_2(G_{n-m}),\Wr)\] to \(\MF_{B_n^2}(\calXr_2(G_n),\Wr)\)
 are isomorphic
$$ \indb_k\circ \mathrm{Id}\times \indb_{m-k}=\indb_m\circ \indb_k\times \mathrm{Id}.$$
\end{corollary}

%%% Local Variables:
%%% mode: latex
%%% TeX-master: "FinalVerMF1_Jan2017"
%%% End:
